\section{Discussion}

We trained recurrent networks made entirely of data-fitted neuron models with a cell-local,
forward-in-time learning rule, and tested every comparison in a protocol locked before the
held-out cohort was used. Four results stand out.

\paragraph{The eligibility trace must contain the fitted adaptation.}
Almost every fitted cell adapts. A local rule whose eligibility trace ignores the cell's own
spike-history filter lost 7--17\pp of out-of-distribution accuracy and fell to the level of a
frozen reservoir. The reason is credit assignment: with $\beta$, the local estimate stays
closely aligned with the exact gradient (cosine up to 0.94 with symmetric feedback), and the
benefit is larger at cells with stronger adaptation loop gain. This extends a point made for
single-threshold adaptive spiking neurons
\citep{bellec2020eprop,bellec2018lsnn,salaj2021sfa}: it holds for heterogeneous, multi-timescale
history filters fitted to real cells. For a synapse, the requirement is concrete. Its trace
must decay with the postsynaptic cell's own adaptation kinetics, not only with membrane or
synaptic time constants.

\paragraph{Where local learning falls short.}
On the two tasks that require combining stored information with a later cue, the local rule
trailed BPTT by 11--22\pp. Symmetric feedback did not close that gap, which points to the paths
through other cells that the rule ignores, rather than to random feedback. The frozen reservoir
was a strong baseline. Population codes show why: the slow fitted filters keep both input
streams decodable through the delay in every network, so learning matters mainly in the response
phase. A frozen reservoir is therefore the right control for claims about recurrent plasticity in
networks of slow cells; in our data, local plasticity beat it significantly in half of the
comparisons.

\paragraph{Biological specificity: a negative result with one exception.}
Once each manipulated bank was recalibrated to the same operating point, the natural pairing of
each cell's parameters and the type-level wiring had no detectable effect. Most such effects
were bounded within $\pm$5\pp. Without recalibration, shuffled networks ran at twice the rate
with about 15 times more saturation. That is a different regime, and earlier versions of this
project had read differences between such networks as effects of cell identity. Heterogeneity
was the exception. Replacing cells by their type means cost BPTT up to 25\pp on the harder tasks,
in line with reports that heterogeneity aids learning \citep{pereznieves2021}, but it did not
affect the local rule. One possible reading is that the diversity of timescales is only useful
to a learner that can assign credit through the network. A local rule that cannot do so does not
benefit from it. Type-mean networks also fired more slowly, so part of this effect may be an
operating-point difference that our calibration did not remove.

\paragraph{Operating point.}
The calibrated baseline rate, a choice that every data-constrained rate model makes, mattered
for the local rule in a task-dependent way and not for BPTT. It should be reported and
varied, not fixed silently.

\paragraph{Limitations.}
\begin{itemize}
\item \emph{The models are rate models.} The GFR models reproduce each cell's firing-rate
response, not spikes, dendrites or conductance-based synapses. The adaptation we find to matter
is the fitted rate-history filter. Spike-level adaptation mechanisms are not modelled.
\item \emph{The cells and the tasks come from different areas.} The cells are mostly from mouse
visual cortex, and the tasks are abstract working-memory tasks. We make no claim that these
circuits solve these tasks in vivo.
\item \emph{The wiring is a prior, not data.} Type-level connection probabilities are modelling
assumptions (Table~\ref{tab:priors}). Measured probabilities were not used because the numbers
we could trace were available only as figures.
\item \emph{The fits have a known split caveat.} Upstream, one stimulus protocol contributes to
both fitting and validation. We selected cells on the released validation criterion and claim no
new patch-clamp validation.
\item \emph{Optimization is a confound we only bound.} Adam with global clipping is not local. We
therefore repeated every learning-rule comparison with per-synapse RMS normalization, which is
local, but no optimizer is a model of biological plasticity. The local rule also uses a
broadcast error and fixed random feedback; we did not test reward-based or neuromodulatory
variants.
\item \emph{Scale.} Networks have 100 cells, with 16 paired seeds per condition. Each network uses
80 of 170 pyramidal cells, so networks share cells: inference is about networks built from this
cohort, not about new cells.
\item \emph{The operating point is a modelling choice.} Calibrating every cell to a fixed fraction
of its own ceiling sets the type-specific rate hierarchy by construction. Family D shows how much
results depend on this choice.
\end{itemize}


\paragraph{Outlook.}
The same pipeline can test spiking versions of the fitted cells, measured connectivity once it
is available in tabular form, reward-modulated rather than supervised signals, and larger
networks. The protocol, code and all run outputs are available for reuse.
